\begin{figure}[htbp]
	\centering
	\begin{tikzpicture}[scale=0.85]
		\node[draw, circle, fill=blue!20, minimum size=0.7cm] (1) at (0,0) {1};
		\node[draw, circle, fill=blue!20, minimum size=0.7cm] (2) at (1.5,0) {2};
		\node[draw, circle, fill=blue!20, minimum size=0.7cm] (3) at (3,0) {3};
		\node[draw, circle, fill=blue!20, minimum size=0.7cm] (4) at (4.5,0) {4};
		\draw[thick] (1) -- (2) -- (3) -- (4);
		\draw[thick] (1) -- (3);
		\draw[thick] (2) -- (4);
		\node[below=0.5cm] at (2.25,-1) {(a) Complete graph $K_4$};

		\node[draw, ellipse, dashed, thick, fill=green!15, minimum width=5cm, minimum height=1.8cm] (h) at (9.5,0) {};
		\node[draw, circle, fill=blue!20, minimum size=0.7cm] (n1) at (7.5,0) {1};
		\node[draw, circle, fill=blue!20, minimum size=0.7cm] (n2) at (8.8,0) {2};
		\node[draw, circle, fill=blue!20, minimum size=0.7cm] (n3) at (10.1,0) {3};
		\node[draw, circle, fill=blue!20, minimum size=0.7cm] (n4) at (11.4,0) {4};
		\node[below=0.5cm] at (9.5,-1) {(b) GHZ hypergraph};

		\node[draw, rounded corners, fill=yellow!20, text width=4cm, align=center] at (6,-4) {
			Single hyperedge captures $|GHZ_N\rangle$\\
			$\frac{1}{\sqrt{2}}(|0\rangle^{\otimes N}+|1\rangle^{\otimes N})$
		};
	\end{tikzpicture}
	\caption{GHZ state representation: (a) Complete graph $K_4$ requires many pairwise edges. (b) GHZ hypergraph represents the same state with a single hyperedge connecting all $N$ qubits, capturing the genuine N-partite entanglement.}
	\label{fig:ghz-hypergraph}
\end{figure}
